\section{Oracularization}
\label{sec:oracle}

\newcommand{\tvora}{\hat{\verifier}^\ora}
\newcommand{\tsora}{\hat{\sampler}^\ora}
\newcommand{\tdora}{\hat{\decider}^\ora}

\subsection{Overview}

In this section we introduce the \emph{oracularization} transformation.
At a high level, the oracularization $\game^\ora$ of a nonlocal game $\game$ is
intended to implement the following: one player (called the \emph{oracle
  player}) is supposed to receive questions $(x,y)$ meant for both players in
the original game $\game$, and the other player (called the \emph{isolated
  player}) only receives either $x$ or $y$ (but not both), along with a label
indicating which player in the original game the question is associated with (we
refer to such players as the \emph{original} players, e.g.\ ``original $\alice$
player'' and ``original $\bob$ player'').
The oracle player is supposed to respond with an answer pair $(a,b)$, and the
other player is supposed to respond with an answer $c$.
The oracle and isolated player win the oracularized game if $(x,y,a,b)$
satisfies the predicate of the original game $\game$ and the isolated player's
answer is consistent with the oracle player's answer.

The oracularization step is needed in preparation to the next section, in which
we perform answer reduction on the introspection game.
To implement answer reduction we need at least one player to be able to compute
a proof, in the form of a PCP, that the decider of the original game would have
accepted the questions $(x,y)$ and answers $(a,b)$.
This requires the player to have access to both questions, and be able to
compute both answers.

\subsection{Oracularizing normal form verifiers}
\label{sec:orac-def}

Let $\verifier = (\sampler,\decider)$ be a normal form verifier such that
$\sampler$ is an $\ell$-level sampler for some $\ell \geq 0$.
We specify the \emph{typed oracularized verifier} $\tvora =
(\tsora,\tdora)$ associated with $\verifier$ as follows. 


\paragraph{Sampler.}
Define the type set $\type^\ora = \{\oracle,\alice,\bob\}$.
(In the remainder of this section we refer to the types in $\type^\ora$ as
\emph{roles}.)
Define the type graph $G^\ora$ that is the complete graph on vertex set
$\type^\ora$ (including self-loops on all vertices).
Define the $\type^\ora$-type sampler $\tsora$ as follows.
For a fixed index $n \in \N$, let $V$ be the ambient space of $\sampler$ and
$L^w$ for $w \in \ab $ be the pair of CL functions of $\sampler$ on index $n$.

Define two $\type^\ora$-typed families of CL functions $\{ L^{w}_\tvar: V \to V
\}$, for $w \in \{\alice,\bob\}$ and $\tvar \in \type^\ora$, as follows:
\begin{equation*}
	L^{w}_\tvar =
  \begin{cases}
		L^\tvar & \text { if $\tvar \in \{\alice,\bob\}$ }, \\
		\mathrm{Id} & \text { if $\tvar = \oracle$ }.
	\end{cases}
\end{equation*}
In other words, if a player gets the type $\tvar \in \{\alice,\bob\}$, then they
get the question that original player $\tvar$ would have received in the game
played by $\verifier_n$.
If they get type $\tvar = \oracle$, then they get the entire seed $z$ that is
used by the sampler $\sampler$, from which they can compute both $L^\alice(z)$
and $L^\bob(z)$, the pair of questions sampled for the players in game
$\verifier_n$.


By definition, the sampler distribution $\mu_{\tsora,\, n}$ has the
following properties.
\begin{enumerate}
\item Conditioned on both players receiving the $\oracle$ role, both players
  receive $z$ for a uniformly random $z \in V$.
		
\item Conditioned on both players receiving the isolated player role, the
  player(s) with role $\alice$ (respectively,~$\bob$) receives $L^\alice(z)$
  (respectively, $L^\bob(z)$) for a uniformly random $z \in V$.
		
\item Conditioned on player $w\in\{\alice,\bob\}$ receiving the $\oracle$ role
  and player $\overline{w}$ receiving the isolated player role, their question
  tuple is distributed according to $\bigl( (\oracle,z), (v,L^v(z)) \bigr)$ if
  $w=\alice$ and $\bigl( (v, L^v(z)), (\oracle,z) \bigr)$ if $w=\bob$, where $z
  \in V$ is uniformly random and $v$ indicates the role of player
  $\overline{w}$.
\end{enumerate}

\paragraph{Decider.}
The typed decider $\tdora$ is specified in Figure~\ref{fig:oracle-decider}.

\begin{figure}[!htb]
  \centering
  \begin{gamespec}
    Input to decider $\tdora$: $(n, \tvar_\alice, x_\alice , \tvar_\bob,
    x_\bob ,a_\alice, a_\bob)$.
    For $w \in \{\alice,\bob\}$, if $\tvar_w = \oracle$, then parse $a_w$ as a
    pair $(a_{w, \alice}, a_{w, \bob})$.
    Perform the following steps sequentially.

    \begin{enumerate}
    \item (\textbf{Game check}).
      \label{enu:oracle-game}
      For all $w \in \{\alice,\bob\}$, if $\tvar_w = \oracle$, then compute
      $x_{w,v} = L^v(x_w)$ for $v \in \{\alice,\bob\}$.
      If $\decider$ rejects $(n,
      x_{w,\alice},x_{w,\bob},a_{w,\alice},a_{w,\bob})$, then reject.

    \item (\textbf{Consistency checks}).
      \label{enu:oracle-full-consistency}
      \begin{enumerate}[nosep]
      \item \label{enu:check-same} If $\tvar_\alice = \tvar_\bob$ and $a_\alice
        \neq a_\bob$, then reject.
      \item \label{enu:oracle-versus-player} If for some $w \in
        \{\alice,\bob\}$, $\tvar_w = \oracle$, $\tvar_{\overline{w}} \in
        \{\alice,\bob\}$, and $a_{w,\, \tvar_{\overline{w}}} \neq
        a_{\overline{w}}$, then reject.
      \end{enumerate}
    \item \label{enu:oracle-acc} Accept if none of the preceding steps rejects.
    \end{enumerate}
  \end{gamespec}
  \caption{Specification of the typed decider $\tdora$.}
  \label{fig:oracle-decider}
\end{figure}

\subsection{Completeness and complexity of the oracularized verifier}

We determine the complexity of the oracularized verifier and
establish the completeness property.

\begin{theorem}[Completeness and complexity of the oracularized verifier]
  \label{thm:oracle-completeness}
  Let $\verifier = (\sampler, \decider)$ be a normal form verifier.
  Let $\tvora = (\tsora,\tdora)$ be the corresponding
  typed oracularized verifier.
  Then the following hold.
  \begin{itemize}
  \item (\textbf{Completeness}) For all $n \in \N$, if $\verifier_n$ has a PCC
    strategy of value $1$, then $\tvora_n$ has an
    SPCC strategy of value $1$.
  
  \item (\textbf{Sampler complexity}) The sampler $\tsora$ depends only
    on $\sampler$ (and not on $\decider$).
    Moreover, the time complexity of $\tsora$ satisfies
    \begin{equation*}
      \TIME_{\tsora}(n) = O \left (\TIME_{\sampler}(n)
      \right),
    \end{equation*}
    Furthermore, if $\sampler$ is an $\ell$-level sampler, then
    $\tsora$ is a $\ell$-level typed sampler.

  \item (\textbf{Decider complexity}) The time complexity of $\tdora$
    satisfies
    \begin{equation*}
      \TIME_{\tdora}(n) = \poly \left (\TIME_{\decider}(n),
        \TIME_\sampler(n) \right)\;.
    \end{equation*}

	\item (\textbf{Efficient computability}) There is a Turing machine
    $\ComputeOracleVerifier$ which takes as input $\verifier =
    (\sampler,\decider)$ and returns $\tvora = ( \tsora,
    \tdora)$ in time $\poly(\abs{\verifier})$.
  \end{itemize}
\end{theorem}


\begin{remark}
Unlike with the Introspection transformation, we do not detype the oracularized verifier $\tvora$; this is because the analysis of the Answer Reduction transformation in \Cref{sec:ans} \emph{directly} reduces to the analysis of the typed oracularized verifier $\tvora$ defined here.
\end{remark}

\begin{proof}
  We analyze the completeness and complexity properties of the typed verifier
  $\tvora$.

  \paragraph{Completeness.}
  For $n \in \N$, let $\strategy = (\ket{\psi},A,B)$ be a PCC strategy for
  $\verifier_n$ with value $1$.
  Consider the following symmetric strategy $\strategy^\ora = (\ket{\psi}, M)$
  for $\tvora_n$.
  Depending on the role received, each player performs the following:
  \begin{enumerate}
	\item Suppose the player receives role $v \in \{\alice,\bob\}$ and question
    $x$.
    Then the player performs the measurement that player $v$ would on question
    $x$ according to strategy $\strategy$ to obtain outcome $a$ (either
    $\{A^{x}_a\}$ or $\{B^x_a\}$, depending on $v$).
    The player replies with $a$.
		\label{enu:oracle-honest-isolated}
		
	\item Suppose the player receives role $v = \oracle$ and question $x$.
    The player first computes $y_w = L^w(x)$ for $w \in \{\alice,\bob\}$ where
    for $w\in \ab$, $L^w$ is the CL functions of $\sampler$ corresponding to
    player $w$.
    Then, the player measures using the POVM $\{ M^{\Oracle,\, x}_{a_\alice,\,
      a_\bob}\}$ where
    \begin{equation}
      \label{eq:oracle-honest}
      M^{\oracle,\, x}_{a_\alice,\, a_\bob} = B^{y_\bob}_{a_\bob}\,
      A^{y_\alice}_{a_\alice}\;.
    \end{equation}
    The projectors $A^{y_\alice}_{a_\alice}$ and $B^{y_\bob}_{a_\bob}$
    commute because $(y_\alice, y_\bob)$ is distributed according to
    $\mu_{\sampler,\, n}$ (over the choice of $x$) and $\strategy$ is a
    commuting strategy for $\verifier_n$.
    Thus $M^{\oracle,\, x}_{a_\alice,\, a_\bob}$ is a projector.
    The player replies with $(a_\alice, a_\bob)$.
		\label{enu:oracle-honest-oracle}
  \end{enumerate}

  The strategy $\strategy^\ora$ is symmetric and projective by construction, and
  consistency follows from the consistency of $\strategy$.
  We now argue that the strategy is commuting and has value $1$ in the game
  $\tvora_n$.
  We consider all possible pairs of roles.

  \begin{enumerate}
  \item ({Oracle, isolated}) Suppose without loss of generality that player
    $w=\alice$ gets the $\oracle$ role and player $\overline{w}=\bob$ gets the
    isolated player $\bob$ role.
    Then player $w$ gets question $x$ and player $\overline{w}$ gets question
    $L^\bob(x)$, where $x$ is uniformly sampled from $V$.
    The oracle player computes $y_v = L^v(x)$ for all $v \in \ab$.
    Notice that $(y_\alice, y_\bob)$ is distributed according to
    $\mu_{\sampler,\, n}$.
    The two players return $\bigl( (a_\alice, a_\bob), a'_\bob \bigr)$ with
    probability
    \begin{align*}
      \bra{\psi} M^{\oracle,\, x}_{a_\alice,\, a_\bob} \otimes
      B^{y_\bob}_{a'_\bob} \ket{\psi}
      & = \bra{\psi} B^{y_\bob}_{a_\bob} A^{y_\alice}_{a_\alice} \otimes
        B^{y_\bob}_{a'_\bob} \ket{\psi}\\
      & = \bra{\psi}  A^{y_\alice}_{a_\alice} \otimes
        B^{y_\bob}_{a_\bob} B^{y_\bob}_{a'_\bob} \ket{\psi}\\
      & = \delta_{a_\bob,\, a'_\bob}\, \bra{\psi} A^{y_\alice}_{a_\alice} \otimes
        B^{y_\bob}_{a_\bob} \ket{\psi}\;,
    \end{align*}
    where the first equality uses the definition of $M^{\oracle,\,
      x}_{a_\alice,\, a_\bob}$ from Eq.~\eqref{eq:oracle-honest}, the second
    equality uses the consistency of $\strategy$, and the third equality uses
    the projectivity of $\strategy$.
    Notice that when $a_\bob = a'_\bob$, this is exactly the probability of
    obtaining answers $(a_\alice,a_\bob)$ when player $\alice$ and player $\bob$
    get question pair $(y_\alice,y_\bob)$ in the game $\verifier_n$.
    Since $\strategy$ is value-$1$, the answers satisfy the decision procedure
    of $\verifier_n$ with probability $1$.
    Thus the oracle's answers pass the ``Game check'' of the oracularized
    decider with certainty, and furthermore the oracle's answers are consistent
    with the isolated player's answers and thus pass the ``Consistency check''
    with certainty as well.
	
    Commutativity of $M^{\Oracle,\, x}_{a_\alice,\, a_\bob}$ and
    $B^{y_\bob}_{a_\bob}$ follows from the commutativity of $\strategy$ for the
    game $\verifier_n$.

  \item ({Both oracle}) If both players get the $\oracle$ role, then both
    players receive the same question $x \in V$.
    Using a similar analysis as for the previous item, the players return the
    same answer pair (thus passing the ``Consistency check'') and pass the
    ``Game check''.
    Both players' measurements commute because they are identical.

  \item ({Both isolated}) Suppose that both players receive the same isolated
    player role (e.g., they both receive the isolated player role $\alice$).
    They then perform the same measurements, which produce the same outcomes due
    to the consistency of the strategy $\strategy$, and thus they pass the
    ``Consistency check''.
    Otherwise, suppose that one player receives the $\alice$ role and the other
    player receives the $\bob$ role.
    Then the decider $\tdora$ automatically accepts.
    Furthermore, their measurements commute because their questions are
    distributed according to $\mu_{\sampler,\, n}$, and $\strategy$ is a
    commuting strategy with respect to $\mu_{\sampler,\, n}$.
  \end{enumerate}

  \paragraph{Complexity.}
  It is clear from the definition that $\tsora$ depends only on
  $\sampler$.
  The time complexity of the sampler $\tsora$ is dominated by those of
  the sampler $\sampler$.
  The complexity of $\tdora$ is dominated by the complexity of both $\sampler$ and $\decider$
  and performing consistency checks.
  The sampler $\tsora$ is a $\max\{\ell,1\}$-level sampler because
  $\sampler$ is an $\ell$-level sampler and the new CL functions for
  $\tvar=\oracle$ are $1$-level.

  \paragraph{Efficient computability.}
  The description of $\tsora$ can be computed, in polynomial time, from
  the description of $\sampler$ alone.
  The description of $\tdora$ can be computed in polynomial time from the
  descriptions of $\sampler$ and $\decider$.
  Moreover, in each case the computation amounts to copying the description of
  $\sampler$ and $\decider$ respectively, and adding constant-sized additional
  instructions.

\end{proof}

\subsection{Soundness of the oracularized verifier}

\begin{theorem}[Soundness of the oracularized verifier]
  \label{thm:oracle-soundness}
  Let $\verifier = (\sampler, \decider)$ be a normal form verifier and
  $\tvora = (\tsora,\tdora)$ be the corresponding typed
  oracularized verifier.
  Then there exists a function $\delta(\eps) = \poly(\eps)$ such that for all $n
  \in \N$ the following hold.
  \begin{enumerate}
  \item If $\val^*(\tvora_n) > 1 - \eps$, then $\val^*(\verifier_n) \geq
    1 - \delta(\eps)$.
	\item For all $\eps > 0$, we have that
    \begin{equation*}
      \Ent \bigl( \tvora_n, 1 - \eps \bigr) \geq
      \Ent(\verifier_{n}, 1 - \delta(\eps))
    \end{equation*}
    where $\Ent(\cdot)$ is as in \cref{def:ent}.
  \end{enumerate}
\end{theorem}

\begin{proof}
	Fix $n \in \N$.
  Let $\strategy^\ora = \bigl( \ket{\psi}, A, B \bigr)$ be a strategy
  for $\tvora_n$ with value $1 - \eps$ for some $0 < \eps \leq 1$. Using Lemma~\ref{lem:symmetric-strat} we may without loss of generality assume that $\strategy^\ora$ is projective. 
  Let $(\tvar,x) \in \type^\ora \times V$ be a question to player $w=\alice$.
  In the event that $\tvar = \oracle$ (which occurs with probability $1/3$), let
  $y_v = L^v(x)$ for each $v \in \{\alice,\bob\}$.
  From the consistency check performed by $\tdora$ and
  item~\ref{item:consistency-implies-approx} of Fact~\ref{fact:agreement}, we
  have that for all $v \in \{\alice,\bob\}$ and on average over $x$ sampled by
  $\tsora$,
  \begin{equation}
    \label{eq:oracle-1}
    A^{\oracle,\, x}_{a_v} \otimes \Id_\bob  \approx_\eps
    \Id_\alice \otimes B^{v,\, y_v}_{a_v}\;.
  \end{equation}     
  Here, we used that with probability $1/9$ player $w=\alice$ gets the $\oracle$
  role and player $\overline{w}=\bob$ gets the isolated player $v$ role;
  conditioned on this, player $\bob$ gets question $y_v$.
	
	Using 
	Fact~\ref{fact:add-a-proj} on~\eqref{eq:oracle-1}, with the $C$ operators from  Fact~\ref{fact:add-a-proj} chosen as $\{A^{\oracle,\, x}_{a_\alice,\,
    a_\bob}\}$ here we get
		  \begin{align}
   A^{\oracle,\, x}_{a_\alice,a_\bob} \otimes \Id_\bob  &=  A^{\oracle,\, x}_{a_\alice,a_\bob} \cdot\big( A^{\oracle,\, x}_{a_\bob} \otimes \Id_\bob \big) \notag\\
	&\approx_\eps
   A^{\oracle,\, x}_{a_\alice,a_\bob}\cdot\big( \Id_\alice \otimes B^{\bob,\, y_\bob}_{a_\bob}\big) \notag\\
	&= A^{\oracle,\, x}_{a_\alice,a_\bob}\otimes B^{\bob,\, y_\bob}_{a_\bob}\;,   \label{eq:oracle-1b}
  \end{align}   
		where for the first equality we used that the POVM elements $\{A^{\oracle,\, x}_{a_\alice,\,
    a_\bob}\}$ are projective. Using a similar calculation for the second and third approximations below we get
  \begin{equation}
    \label{eq:oracle-2}
    \begin{split}
      A^{\oracle,\, x}_{a_\alice,\, a_\bob} \otimes \Id_\bob
      & \approx_{\eps} A^{\oracle,\, x}_{a_\alice,\, a_\bob}  \otimes
      B^{\bob,\, y_\bob}_{a_\bob} \\
      &\approx_{\eps} A^{\oracle,\, x}_{a_\alice} \otimes
      B^{\bob,\, y_\bob}_{a_\bob} \\
      & \approx_\eps \Id_\alice \otimes
      B^{\bob,\, y_\bob}_{a_\bob} B^{\alice,\, y_\alice}_{a_\alice} \\
      & \approx_\eps \Id_\alice \otimes
      B^{\alice,\, y_\alice}_{a_\alice} B^{\bob,\, y_\bob}_{a_\bob}\;,
    \end{split}
  \end{equation}
	where the last line uses that starting from $A^{\oracle,\, x}_{a_\alice,\, a_\bob} \otimes \Id_\bob$ we may perform all operations leading to the before-last line hwile reversing the order in which $B^\bob$ and $B^\alice$ are introduced. 
	
  Using \cref{enu:check-same} of the consistency check, we have that on
  average over a random $y = L^\alice(x) \in L^\alice(V)$,
  \begin{equation}
    \label{eq:oracle-3}
    A^{\alice,\, y}_{a} \otimes I_\bob \approx_\eps
    \Id_\alice \otimes B^{\alice,\, y}_{a}\;.
  \end{equation}
  Define, for all $y \in L^\alice(V)$, measurement operators $\{C^y_a\}_a$ where
  $C^y_a = A^{\alice,\, y}_a$.
  Similarly, define $D^y_a = B^{\bob,\, y}_a$.
  This defines a strategy $\strategy = (\ket{\psi},C,D)$ for the game
  $\verifier_n$ that we now argue succeeds with high probability.
    
  Let $x \in V$ be uniformly random.
  Let $y_\alice = L^\alice(x)$ and $y_\bob =L^\bob(x)$.
  \begin{align*}
    A^{\oracle,\, x}_{a_\alice,\, a_\bob} \otimes \Id_\bob
    & \approx_\eps \Id_\alice \otimes
      B^{\bob,\, y_\bob}_{a_\bob} B^{\alice,\, y_\alice}_{a_\alice} \\
    & \approx_{\eps} A^{\alice,\, y_\alice}_{a_\alice} \otimes
      B^{\bob,\, y_\bob}_{a_\bob} \\
    & = C^{y_\alice}_{a_\alice} \otimes D^{y_\bob}_{a_\bob}\;.
  \end{align*}
	The first approximation follows from Equation~\eqref{eq:oracle-2}.
  The second approximation follows from Equation~\eqref{eq:oracle-3} and
  Fact~\ref{fact:add-a-proj} (where we let $C^y_b$ in Fact~\ref{fact:add-a-proj}
  represent $\Id\ot B^{\bob,\, y_\bob}_{a_\bob}$).
  The last equality follows from definition of $C^{y_\alice}_{a_\alice}$ and
  $D^{y_\bob}_{a_\bob}$.
  The pair of questions $(y_\alice, y_\bob) \in V \times V$ is distributed
  according to $\mu_{\sampler,\, n}$.

  The game check part of $\tdora$ succeeds with probability $1 -
  O(\eps)$, which implies that the answer pair $(a_\alice, a_\bob)$ that arises
  from the measurement $A^{\oracle,\, x}_{a_\alice,\, a_\bob} \otimes \Id_\bob$
  is accepted by the decider $\decider$ on question pair $(y_\alice,y_\bob)$
  with probability $1 - O(\eps)$.
  This in turn implies that the strategy $\strategy = (\ket{\psi},C,D)$ succeeds
  with probability $1 - O(\sqrt{\eps})$ in the game $\verifier_n$.
  As an additional consequence, the Schmidt rank of $\ket{\psi}$ must be at
  least $\Ent(\verifier_n, 1 - O(\sqrt{\eps}))$.
\end{proof}

 
\section{Answer Reduction}
